\documentclass[11pt,a4paper]{article}

\usepackage[margin=1.05in]{geometry}
\usepackage{amsmath,amssymb,amsthm}
\usepackage{mathtools}
\usepackage{stmaryrd}
\usepackage{booktabs}
\usepackage{array}
\usepackage{listings}
\usepackage{xcolor}
\usepackage{graphicx}
\usepackage[colorlinks=true,linkcolor=black,citecolor=black,urlcolor=blue]{hyperref}
\usepackage{enumitem}
\usepackage[T1]{fontenc}
\usepackage[protrusion=true,expansion=false]{microtype}
\setlength{\emergencystretch}{2.5em}

%%%%%%%%%%%%%%%%%%%%%%%%%%%%%%%%%%%%%%%%%%%%%%%%%%%%%%%%%%%%%%%%%%%%%%%%%%%%%%%
% Theorem environments
%%%%%%%%%%%%%%%%%%%%%%%%%%%%%%%%%%%%%%%%%%%%%%%%%%%%%%%%%%%%%%%%%%%%%%%%%%%%%%%
\theoremstyle{plain}
\newtheorem{theorem}{Theorem}[section]
\newtheorem{proposition}[theorem]{Proposition}
\newtheorem{lemma}[theorem]{Lemma}
\newtheorem{corollary}[theorem]{Corollary}
\newtheorem{conjecture}[theorem]{Conjecture}

\theoremstyle{definition}
\newtheorem{definition}[theorem]{Definition}
\newtheorem{example}[theorem]{Example}
\newtheorem{principle}[theorem]{Principle}
\newtheorem{obligation}[theorem]{Obligation}
\newtheorem{observation}[theorem]{Observation}

\theoremstyle{remark}
\newtheorem{remark}[theorem]{Remark}

%%%%%%%%%%%%%%%%%%%%%%%%%%%%%%%%%%%%%%%%%%%%%%%%%%%%%%%%%%%%%%%%%%%%%%%%%%%%%%%
% Macros.  NOTATION RULES OBSERVED THROUGHOUT:
%   * the rho calculus is never written with the Greek letter rho
%   * quoting is @P, dropping is *x; corner quotes are pre-2016 notation
%%%%%%%%%%%%%%%%%%%%%%%%%%%%%%%%%%%%%%%%%%%%%%%%%%%%%%%%%%%%%%%%%%%%%%%%%%%%%%%
\newcommand{\rhoc}{\textsf{rho}}
\newcommand{\sem}[1]{\llbracket #1 \rrbracket}
\newcommand{\V}{\mathbb{V}}
\newcommand{\Bool}{\mathbb{B}}
\newcommand{\Rnn}{\mathbb{R}_{\geq 0}}
\newcommand{\Cx}{\mathbb{C}}
\newcommand{\Nat}{\mathbb{N}}
\newcommand{\Proc}{\mathcal{P}}
\newcommand{\Names}{\mathcal{N}}
\newcommand{\Reach}{\mathrm{Reach}}
\newcommand{\supp}{\mathrm{supp}}
\newcommand{\scong}{\equiv}
\newcommand{\tensorv}{\otimes}
\newcommand{\sumv}{\oplus}
\newcommand{\Cand}{\mathrm{Cand}}
\newcommand{\Out}{\mathrm{out}}
\newcommand{\crisp}[1]{\lceil #1 \rceil}
\newcommand{\dok}{\textsf{do}}
\newcommand{\ink}{\textsf{in}}
\newcommand{\fork}{\textsf{for}}
\newcommand{\where}{\textsf{where}}
\newcommand{\PA}{\mathrm{PA}}
\newcommand{\dsep}{\perp\!\!\!\perp}

\newcommand{\lolli}{\mathrel{\reflectbox{$\multimap$}}}

\lstdefinestyle{rholang}{
  basicstyle=\ttfamily\small,
  keywordstyle=\bfseries,
  morekeywords={for,new,in,where,contract,match,if,else,do},
  columns=fullflexible,
  keepspaces=true,
  showstringspaces=false,
  breaklines=true,
  frame=single,
  framesep=5pt,
  rulecolor=\color{black!35},
  xleftmargin=6pt,
  aboveskip=9pt,
  belowskip=9pt
}
\lstset{style=rholang}

%%%%%%%%%%%%%%%%%%%%%%%%%%%%%%%%%%%%%%%%%%%%%%%%%%%%%%%%%%%%%%%%%%%%%%%%%%%%%%%

\title{\bfseries Interventions at the Cut\\[4pt]
\large Pearl's \textsf{do} as a scoped guard in the rho calculus with graded
where-clauses: an exact embedding of causal Bayesian networks, its
$\sigma$-calculus generalisation, and the fragments that make it run}

\author{L.\,G.\ Meredith\\ \small F1R3FLY.io}

\date{F1R3FLY.io working note --- November 2026\\[2pt]
\normalsize Draft 1}

\begin{document}
\maketitle

\begin{abstract}
A program in the rho calculus with graded where-clauses, valued over the
non-negative reals, already carries a causal model: the likelihood of a causal
history is a product over firing sites of a locally normalised factor, each
site's factor is its clause, each site's parents are the data its pattern and
clause read, and the resolver that settles its race is the exogenous noise.
This note makes the model explicit and adds the one thing Pearl's calculus has
that the clause language lacked --- an operator that intervenes.

The operator is a scoped decoration on parallel composition,
$\dok(\,y_1 \models \varphi_1 \ \&\ \cdots\ \&\ y_n \models \varphi_n\,)\
\ink\ P$, read as: whenever a race for the binding of $y_i$ arises inside
$P$, resolve it among the candidates satisfying $\varphi_i$, with the site's own
weights restricted to that set. It is placed at the cut because an aperture is a
property of the cut and of neither agent flanking it, and because whether a
particular cut will ever occur is undecidable even in the sequential fragment,
so an intervention cannot be a surgery on a known node and must be a policy
conditional on an event. The construct lowers into the where-clause by the
normal-form theorem for weight maps, so it is a conservative extension, and it
sits beside the cost membrane as a second index on a polymorphic cut.

With it, Pearl's causal Bayesian networks over finite domains embed exactly: a
variable is a join on its parents' channels together with an ancilla race over
its domain, the conditional probability table is the clause in disjunctive
normal form, the term is staged so that each gate is its own contention set,
and the truncated factorisation is the likelihood of a causal history.
Intervening on the ancilla binding deletes the incoming edges on the nose:
normalisation sends the one surviving weight to $1$ for every value of the
parents, and the positivity condition Pearl needs appears as a receipt that
can starve. The three rules of the do-calculus are then theorems about the
embedded fragment rather than axioms, and counterfactuals are replay from a
recorded pre-state under the intervention. The construction is the one-qubit
gadget of the quantum paper with real weights, and we say exactly what the
resemblance amounts to. Restricting to a set rather than a point, or supplying
a graded constraint, lands in Correa and Bareinboim's $\sigma$-calculus, whose
induced graph we read off the free names of the constraint.

The naive interpretation of the construct is expensive in three independent
ways --- the constraint language can contain modalities and fixed points, the
scope can be dynamic through reflection, and restriction changes the
normalisation of contention sets whose enumeration is permanent-shaped --- and
whether a given $\dok$ is vacuous is a reachability question. We propose
fragments of the constraint language graded by evaluation cost, and type-based
restrictions on the scope (staging, flatness, finite-domain channel types,
namespace confinement) under which lowering is exact, checks are constant-time,
and the semantics is Pearl's. We close by proposing the mechanism as the
constructive content owed by the companion note on inhabiting choice: rung two
of Pearl's ladder is a graded race under a $\dok$, rung three is its replay, and
the vending-machine refusal that no trace can reveal is a receipt that starves
under a one-line intervention.
\end{abstract}

\tableofcontents

%%%%%%%%%%%%%%%%%%%%%%%%%%%%%%%%%%%%%%%%%%%%%%%%%%%%%%%%%%%%%%%%%%%%%%%%%%%%%%%
\section{Introduction}\label{sec:intro}

\subsection{Where the choice is}

A \rhoc{} program does not determine its own execution. Write
\[
  x!(Q_1) \mid \fork(y \leftarrow x)P \mid x!(Q_2)
\]
and the operational semantics permits two reductions without saying which
happens or with what tendency. The graded where-clause programme
\cite{GradedV4,GradedQuantum,GradedPLN} takes this silence as a parameter: the
\where{} clause of the for-comprehension is valued in a commutative semiring and
evaluated once per candidate match, so that the programmer may write the
resolution mechanism down rather than inherit it. Over the Booleans this is
rholang; over $\Cx$ it is a syntax for arranging interference; over $\Rnn$ it is
stochastic execution, and the companion paper \cite{GradedPLN} showed that the
real instance is the one in which a learner lives, because it is the one with a
likelihood to differentiate.

The present note observes that the real instance is also the one in which a
\emph{causal} model lives, in the precise sense of Pearl
\cite{Pearl95,Pearl09}, and that it has been living there since the clause was
introduced. What was missing was not a model but an operator: the calculus
could observe its races and weight them, but had no syntax for \emph{forcing}
one. Pearl's $\dok$ is that syntax, and this note gives it a home.

\subsection{The observation}

The likelihood of a causal history in \cite{GradedPLN} is
\begin{equation}\label{eq:lik}
  \Pr\nolimits_\theta(h) \;=\; \prod_{i=1}^{\ell}
  \frac{w(\theta,\sigma_i)}{\sum_{\tau \in \Gamma_i} w(\theta,\tau)} ,
\end{equation}
a product over firing sites of a factor normalised within the site's contention
set $\Gamma_i$. Read with Pearl's dictionary this is a Bayesian-network
factorisation. The variable is the site $r$ and its value is the matching
$\sigma$ that fired. Its parents are the contents of $\Gamma_r$ --- what rests
on the channels the pattern and the clause read. Its mechanism is the clause
$\psi_r$. Its exogenous noise is the resolver's draw. And the thesis of the
whole programme, that the resolution mechanism is a parameter supplied from
outside the term, is Pearl's exogeneity: $U$ is what the model does not
explain.

Three consequences follow before any construction, and we state them here
because they govern what follows.

\begin{itemize}[leftmargin=1.4em,itemsep=3pt]
\item \emph{The parent set is syntactic.} Pearl must posit the graph. Here the
parents of a site are the names its pattern and the crisp subformulae of its
clause read; the graph is written in the program. The one construct that can
smuggle in an unlisted parent is reflection --- $*y$ dropping received code ---
which is the same suspect the choice-formulae note \cite{ChoiceFormulae}
identified for subject reduction.
\item \emph{Observation and intervention are already priced differently.}
Conditioning on ``the race at $r$ went $x$'' keeps the histories in which it did
and renormalises, which \cite{GradedV4} showed is post-selection and charged at
$1/\Pr(x)$ retries. Intervention replaces a column block of the transfer
operator and renormalises nothing. The two differ exactly when histories
producing $x$ also influence the outcome by another route, which is a
confounder; the paradigm confounder in the programme's own examples is a belief
datum read by two sites.
\item \emph{The do-calculus is Paper II territory.} Pearl's rules presuppose a
measure over histories. Over $\Cx$ a reachable outcome may carry amplitude zero,
so there is no per-branch probability \cite[Principle~3.5]{GradedPLN}, and the
counterpart on the complex side is the quantum-causal-models literature
\cite{CostaShrapnel,AllenEtAl}, a different object. Everything below is over
$\Rnn$.
\end{itemize}

\subsection{Contributions}

\begin{enumerate}[leftmargin=1.4em,itemsep=3pt]
\item \textbf{Interventions have one carrier} (\S\ref{sec:carrier}). An
intervention is a change of section over the candidate fibre of a site. The
fibre at a deterministic site is a point, so deterministic communication is
intervention-transparent (the dual of amplitude-transparency,
\cite[Prop.~8.5]{GradedV4}); a context cannot force a race it does not own over
$\Rnn$ (the converse of the aperture principle), and can exactly in the
idempotent algebras; and the observation/intervention gap --- the whole content
of Pearl's calculus --- lives only on races.
\item \textbf{The construct} (\S\ref{sec:construct}). A scoped decoration on
the cut, $\dok(C)\ \ink\ P$, with a membrane semantics and a lowering into the
where-clause that agree on the reflection-free scope. It is a conservative
extension by the normal-form theorem for weight maps, it composes by
conjunction, and it is the operator that moves a formula from specification
position to guard position in the sense of \cite{Omnibus}.
\item \textbf{The embedding} (\S\ref{sec:pearl}). Causal Bayesian networks over
finite domains are the $\Rnn$ instance of the graded clause on staged terms;
$\dok$ on the ancilla binding is graph surgery on the nose
(Theorem~\ref{thm:surgery}); the three rules are derived
(Corollary~\ref{cor:rules}); counterfactuals are replay
(Theorem~\ref{thm:cf}). Positivity appears as starvation.
\item \textbf{The Shor gadget} (Remark~\ref{rem:shor}). The variable gadget is
the one-qubit gadget of \cite{GradedQuantum} with a stochastic matrix in place
of a unitary one, and the resemblance is a theorem about the semiring rather
than an analogy.
\item \textbf{The $\sigma$-calculus} (\S\ref{sec:sigma}). Set-valued and graded
constraints are Correa--Bareinboim $\sigma$-interventions
\cite{CorreaBareinboim}; the induced graph is read off the free names of the
constraint; Pearl's atomic $\dok$ is the singleton case.
\item \textbf{Complexity, and the fragments that tame it}
(\S\ref{sec:complexity}). The naive interpretation pays for modalities and
fixed points in the constraint, for dynamic scope through reflection, and for
permanent-shaped renormalisation of contention sets; vacuity of a $\dok$ is a
reachability question. We grade the constraint language into four fragments by
evaluation cost and impose four type disciplines on the scope under which
lowering is exact and the semantics is Pearl's
(Proposition~\ref{prop:practical}).
\item \textbf{Execution} (\S\ref{sec:execution}). The construct lowers in the
elaboration pass the node already has, its membrane form is the seal tree's
twin, an intervention is a COMM and is priced as one, and a counterfactual is
the trace-replay the validators already perform.
\item \textbf{Inhabiting choice, discharged} (\S\ref{sec:inhabiting}). The
companion note \cite{Inhabiting} aligned Pearl's ladder with the
linear/branching/causal spectrum by argument; this note supplies the
constructive content, and the vending-machine benchmark becomes a one-line
program.
\end{enumerate}

\subsection{What this note does not do}

It does not treat continuous domains, which need the DDL's numeric types to
carry an ancilla as a density. It does not prove the $\sigma$-calculus rules
beyond the staged fragment, where the graph is dynamic and the right objects
are the branching cells of \cite{AbbesBenveniste}; that is stated as an
obligation. It does not fix the token charge of a $\dok$ against the ledger. And
it does not treat the complex instance.

%%%%%%%%%%%%%%%%%%%%%%%%%%%%%%%%%%%%%%%%%%%%%%%%%%%%%%%%%%%%%%%%%%%%%%%%%%%%%%%
\section{The calculus}\label{sec:calculus}

The presentation is self-contained and follows \cite{GradedPLN}; a reader of
that paper may skip to \S\ref{sec:carrier}.

\begin{definition}[Syntax]\label{def:syntax}
Processes $P$ and names $x$ are given by
\[
\begin{array}{rcl}
P &::=& 0 \;\mid\; x!(P) \;\mid\; \fork(\,\vec{y} \leftarrow \vec{x}
  \;\where\; \psi\,)\,P \;\mid\; P \mid P \;\mid\; *x
  \;\mid\; \textsf{new}\ x\ \ink\ P \\[2pt]
x &::=& @P
\end{array}
\]
A receipt $\fork(y_1 \leftarrow x_1 \,\&\, \cdots \,\&\, y_m \leftarrow x_m
\;\where\; \psi)P$ is a \emph{join}: it fires only when every pattern is matched.
The clause $\psi$ is drawn from the graded sort below and takes values in $\V$;
its absence abbreviates $\crisp{\top}$. Receipts are linear by default; the
persistent form $\Leftarrow$ survives its firing.
\end{definition}

\begin{definition}[Structural congruence]
$\scong$ is the least congruence making $(\,\mid\,, 0)$ a commutative monoid and
satisfying $*@P \scong P$ and $@*x = x$, closed under $\alpha$-equivalence and
scope extrusion for \textsf{new}.
\end{definition}

\begin{definition}[Formula sorts]\label{def:sorts}
The \emph{crisp} sort retains the full propositional structure of the
generated logic, including negation and fixed points:
\[
\beta ::= \top \;\mid\; \Out(n)\beta \;\mid\; \mathrm{in}(n)\beta
      \;\mid\; \beta \mid \beta \;\mid\; @\beta
      \;\mid\; \beta \wedge \beta \;\mid\; \neg\beta
      \;\mid\; \langle K\rangle\beta \;\mid\; \nu X.\beta .
\]
The \emph{graded} sort has no negation and takes values in $\V$:
\[
\psi ::= \crisp{\beta} \;\mid\; v \;\mid\; \psi \tensorv \psi
      \;\mid\; \psi \sumv \psi \;\mid\; \langle K\rangle_{\V}\psi
      \;\mid\; X \;\mid\; \mu X.\psi \;\mid\; \nu X.\psi ,
\]
with $v$ ranging over constants of $\V$ and fixed-point variables positive and
guarded. $\crisp{-}$ sends a crisp formula to $1$ where it holds and to $0$
where it does not.
\end{definition}

\begin{definition}[Candidates and contention]\label{def:cand}
Fix a term $P$ and a receipt occurrence $r$ with patterns $y_1,\dots,y_m$ on
channels $x_1,\dots,x_m$. A \emph{candidate} at $r$ is an injective assignment
$\sigma$ of the patterns to messages resting on the corresponding channels;
$\Cand(P,r)$ is the set of candidates. Two candidates \emph{contend} if they
share a message or share a linear receipt; a \emph{contention set} is a
connected component of the contention relation, and it is the unit at which
weights are normalised.
\end{definition}

\begin{definition}[Weighted communication]\label{def:comm}
Write $\psi\sem{\sigma} \in \V$ for the value of the clause with the pattern
variables bound as $\sigma$ prescribes. If $\psi\sem{\sigma} \neq 0$ then
$P \longrightarrow_{\psi\sem{\sigma}} P'$, where $P'$ is the residue of consuming
the assigned messages and $r$ and substituting. A candidate of value $0$ offers
no transition.
\end{definition}

\begin{definition}[Completion over $\Rnn$]\label{def:completion}
Throughout, $\V = (\Rnn,+,\times,0,1)$. Resolution normalises within a
contention set: the probability that a given maximal matching fires is its
weight divided by the total over the set. A contention set is \emph{complete}
when that total is $1$.
\end{definition}

\begin{definition}[Causal history]\label{def:history}
A causal history $h$ of length $\ell$ is a sequence of contention sets and
maximal matchings $(\Gamma_1,\sigma_1),\dots,(\Gamma_\ell,\sigma_\ell)$,
quotiented by the interleaving of independent firings. Its likelihood is
\eqref{eq:lik}.
\end{definition}

The quotient is load-bearing: without it $m$ independent firings contribute the
$m!$ linearisations of a partial order, measured as $2, 6, 24$ in the companion
papers. With it the likelihood is a product over \emph{events}, and an event
structure's configurations are partial orders, hence acyclic. This is the
first point at which the dictionary of \S\ref{sec:intro} must be read with
care: the \emph{event-level} graph is always a DAG; the \emph{variable-level}
graph, which identifies events of a kind, may have cycles when the term
recurses. Pearl applies verbatim at the event level. The variable level is
where the fixed point of \cite{GradedPLN} lives, and where this note imposes
types.

\begin{proposition}[Support adequacy is exact over $\Rnn$ {\cite[Prop.~3.3]{GradedPLN}}]
\label{prop:adequacy}
With all clause constants strictly positive, $\supp\sem{P} = \Reach(P)$: every
reachable outcome carries positive weight, so every branch of every contention
set has a probability.
\end{proposition}

%%%%%%%%%%%%%%%%%%%%%%%%%%%%%%%%%%%%%%%%%%%%%%%%%%%%%%%%%%%%%%%%%%%%%%%%%%%%%%%
\section{Interventions have one carrier}\label{sec:carrier}

Before adding syntax it is worth settling what an intervention can touch, since
the answer constrains the syntax.

\begin{definition}[Intervention, semantically]\label{def:intervention}
An \emph{intervention} at a site $r$ of $P$ is a replacement of the map
$\Cand(P,r) \to \V$ denoted by $\psi_r$ with another such map. It is
\emph{hard} when the replacement is the indicator of a single candidate,
\emph{set-valued} when it is the indicator of a subset, and \emph{soft}
otherwise.
\end{definition}

\begin{proposition}[Deterministic communication is intervention-transparent]
\label{prop:transparent}
At a site with exactly one candidate, every intervention is either the identity
on the normalised dynamics or disables the site. In particular no intervention
at such a site changes which outcome is reached when the site fires.
\end{proposition}

\begin{proof}
The fibre is a singleton; a map on a singleton is determined by whether its
value is zero. Normalisation within the contention set sends any nonzero value
to $1$.
\end{proof}

This is the dual of amplitude-transparency \cite[Prop.~8.5]{GradedV4}: a
deterministic step neither carries weight nor admits a choice. To intervene at
a deterministic site one must change the fibre --- inject a datum or delete one
--- and the moment a datum is injected the site has two candidates and is a
race. Deletion is the degenerate case, resolution toward the empty section.

\begin{proposition}[No hard intervention by context over $\Rnn$]
\label{prop:nocontext}
Let $r$ be a race in $P$. No \rhoc{} context $C[-]$ makes $C[P]$ resolve $r$
to a chosen candidate with probability $1$ under the normalised $\Rnn$
semantics, unless $C$ consumes every competing candidate before $r$ can fire
--- which is itself a race $C$ wins with probability strictly less than $1$. Over
an idempotent resolution algebra (Viterbi, tropical) a context whose eater
carries strictly dominating weight does force the outcome.
\end{proposition}

\begin{proof}
The first claim is the aperture principle \cite[Prin.~2.1]{GradedV4} read in
the forcing direction: a context can add contention and consume contenders, and
consuming a contender is a candidate at a site in contention with $r$. Over
$\Rnn$ every candidate of positive weight has positive probability
(Proposition~\ref{prop:adequacy}), so $r$ fires first with positive probability.
Over an idempotent algebra resolution is argmax, which is attained by a
definable section when the maximum is unique \cite[Prop.~3.4(i)]{ChoiceFormulae}.
\end{proof}

So the carrier of intervention is the candidate set at a receipt, a nontrivial
intervention exists exactly where that set has or is given more than one
element, and over $\Rnn$ hard intervention requires access to the clause ---
white-box access, in the vocabulary of \cite{ExpRegimes}. Two qualifications
belong here. First, it is the \emph{gap} between observation and intervention
that lives only on races: one may replace $x!(1)$ by $x!(2)$ in a confluent
program and change its output with no race anywhere, but that is running a
different term, $\Pr(y\mid x)$ and $\Pr(y\mid \dok(x))$ coincide trivially
because there are no histories to reweight, and none of Pearl's rules has
content there. Second, the clause controls the jump chain and not the holding
times \cite[Rem.~3.2]{GradedPLN}, so an intervention on \emph{when} a site fires
is outside the clause and outside this note unless time is reified as data, at
which point it is a race again.

\begin{principle}[Interventions are metered because they are races]
\label{prin:metered}
A race is a COMM and a COMM costs one unit of phlo. Reading the trace is free;
intervening is charged. This is not a policy choice: intervention has no
carrier other than the event the ledger already prices.
\end{principle}

%%%%%%%%%%%%%%%%%%%%%%%%%%%%%%%%%%%%%%%%%%%%%%%%%%%%%%%%%%%%%%%%%%%%%%%%%%%%%%%
\section{The construct: a decorated cut}\label{sec:construct}

\subsection{Why the cut}

There is always a node in the syntax tree above a race, and it is a parallel
composition: a race needs a receiver and at least two things it could meet. The
graded programme already located the aperture there --- an aperture is a
property of the cut and not of either agent flanking it
\cite[Rem.~2.3]{GradedV4} --- and the cost-accounting programme already made the
cut polymorphic, indexing it by signatures and seals \cite{CostMonad}. An
intervention is a second index on the same operator.

The reason it cannot be anything more specific is a halting argument. Whether
a particular cut will ever occur is undecidable, already in the sequential
fragment, since \rhoc{} is Turing complete and ``this receipt fires'' is a
nontrivial reachability property. So an intervention cannot be a surgery on a
node known in advance; it must be a \emph{policy conditional on an event}:
\emph{if} a race for this binding arises, \emph{then} resolve it thus. Vacuous
if the race never arises. Pearl's nearest object is a dynamic treatment regime
rather than an atomic $\dok$, and the formal home is the unfolding rather than
the DAG.

\subsection{Syntax}

\begin{definition}[The $\dok$ membrane]\label{def:do}
Extend Definition~\ref{def:syntax} with
\[
  P ::= \cdots \;\mid\; \dok(\,C\,)\ \ink\ P, \qquad
  C ::= y \models \varphi \;\mid\; y \models \psi \;\mid\; C \,\&\, C ,
\]
where $y$ is a pattern variable, $\varphi$ a crisp formula and $\psi$ a graded
formula, each with $y$ among its free names. We write $\dok(y = Q)$ for
$\dok(y \models \crisp{y = @Q})$ and $\dok(y \in S)$ for the disjunction over a
finite set.
\end{definition}

\begin{lstlisting}
do( y = Q2 ) in { x!( Q1 ) | for( y <- x )P | x!( Q2 ) }

do( y1 |= phi1 & y2 |= phi2 ) in { P1 | P2 | P3 }
\end{lstlisting}

\subsection{Semantics, twice}

\begin{definition}[Membrane semantics]\label{def:membrane}
Let $r$ be a receipt occurrence inside the scope of $\dok(C)\ \ink\ P$ binding a
variable $y$ constrained in $C$ by $\varphi$ (resp.\ $\psi$). At $r$ the
clause is read as
\[
  \psi_r \tensorv \crisp{\varphi} \qquad\text{(resp. } \psi_r \tensorv \psi
  \text{)},
\]
with the constraint evaluated under the candidate's substitution, and
normalisation within the contention set is taken \emph{after} the restriction.
Reduction is $\dok(C)\ \ink\ P \longrightarrow_w \dok(C)\ \ink\ P'$ whenever
$P \longrightarrow_w P'$ under clauses so read. A receipt in scope binding an
unconstrained variable is unaffected; a receipt arriving in scope by reflection
is constrained like any other.
\end{definition}

\begin{definition}[Lowering]\label{def:lowering}
On a scope $P$ in which every receipt binding a constrained variable is
syntactically present and its channel is a literal or \textsf{new}-bound name,
define $\lfloor \dok(C)\ \ink\ P \rfloor$ to be $P$ with each such receipt's
clause replaced by $\psi_r \tensorv \crisp{\varphi}$ (resp.\ $\psi_r \tensorv
\psi$), and the $\dok$ erased.
\end{definition}

\begin{proposition}[Lowering is exact on flat scopes]\label{prop:lowering}
If $P$ contains no drop $*z$ at a position that can come to contain a receipt,
then $\dok(C)\ \ink\ P$ and $\lfloor \dok(C)\ \ink\ P \rfloor$ have the same
weighted transition system, hence the same graded denotation and the same
likelihood on every causal history.
\end{proposition}

\begin{proof}
Under the hypothesis every receipt that will ever be in scope is present in $P$
up to substitution of names for pattern variables, which does not create or
destroy binders. Definition~\ref{def:membrane} reads each such clause exactly as
Definition~\ref{def:lowering} rewrites it, and the contention relation is
unchanged because restriction removes candidates without adding any.
\end{proof}

The hypothesis is the \emph{flat} fragment of \cite{ExpRegimes}, and the gap
between the two semantics is exactly the dynamic scope that reflection
introduces. We take the membrane as the definition and the lowering as the
optimisation, for the same reason \cite{ChoiceFormulae} stratifies \where{}:
received code is the one route by which an unlisted binder --- and, in the
causal reading, an unlisted parent --- enters.

\begin{remark}[Hygiene]
$\dok(y \ldots)\ \ink\ P$ refers to a name bound inside $P$, so $y$ ceases to be
$\alpha$-convertible independently of the $\dok$. We read the construct as part
of the binding scope: the $\dok$ and every receipt it reaches are renamed
together. The alternative --- labelled receipts addressed by label --- is better
theory and worse ergonomics; the semantics is the same.
\end{remark}

\subsection{Restriction, override, and the normal form}

\begin{observation}[Restrict, not override]\label{obs:restrict}
Definition~\ref{def:membrane} \emph{multiplies} the clause by the constraint.
It does not replace it. For $\fork(y \leftarrow x \;\where\; \crisp{*y > 5})P$
under $\dok(y = 3)$, the restricted clause is $0$ on every candidate and the site
waits; an override would force a $3$. Pearl's atomic $\dok$ severs the mechanism
and so is an override in spirit --- but \S\ref{sec:pearl} shows that on the
embedded models the two coincide, because the embedding puts the mechanism in
the \emph{weights} and the restriction is on an ancilla whose own gate is
$\top$. Where they differ, restriction is the construct of this note and
override is a $\sigma$-intervention (\S\ref{sec:sigma}).
\end{observation}

\begin{proposition}[A compressed $\dok$ is a weight map]\label{prop:dnf}
A constraint $y \models \bigoplus_i v_i \tensorv \crisp{\chi_i}$ with the
$\chi_i$ pairwise exclusive denotes the weight map $(\chi_i, v_i)_i$ of the
weighted-GSLT programme restricted to the fibre at $y$'s site, and every weight
map on that fibre is so denoted. In particular $\dok(y = Q_1 \vee y = Q_2)$ is
the weight map assigning $1$ to two candidates and $0$ to the rest.
\end{proposition}

\begin{proof}
This is \cite[Thm.~7.2]{GradedV4} applied to the constraint rather than the
clause.
\end{proof}

So several interventions compress into one expression for the reason weight
maps have a normal form: a formula has one value. The $\dok$ is a scoped weight
map placed at the cut, and the normal-form theorem is what makes it a
conservative extension.

\begin{proposition}[Composition]\label{prop:compose}
$\dok(C_1)\ \ink\ \dok(C_2)\ \ink\ P$ and $\dok(C_1 \,\&\, C_2)\ \ink\ P$ have the
same weighted transition system; $\&$ is commutative and associative; and
$\dok(y \models \top)\ \ink\ P$ has the transition system of $P$.
\end{proposition}

\begin{proof}
Multiplication in $\Rnn$ is commutative and associative with unit $1$, and
$\crisp{\top} = 1$.
\end{proof}

\begin{remark}[Specification into guard]\label{rem:specguard}
The omnibus \cite[Part~III]{Omnibus} places one formula in four positions:
specification, test, guard, hypothesis. A $\dok(y \models \varphi)$ takes the
$\nu$-property ``every binding of $y$ satisfies $\varphi$'' --- a specification
one would check --- and installs it in guard position at every site in scope,
including sites not yet visible. The $\dok$ is the operator that turns a
specification into a guard, and the Principle that the guard is where
resolution becomes legible is why that is the right place for an intervention
to live.
\end{remark}

%%%%%%%%%%%%%%%%%%%%%%%%%%%%%%%%%%%%%%%%%%%%%%%%%%%%%%%%%%%%%%%%%%%%%%%%%%%%%%%
\section{The embedding of Pearl's calculus}\label{sec:pearl}

\subsection{Causal Bayesian networks}

\begin{definition}[Causal Bayesian network]\label{def:cbn}
A \emph{causal Bayesian network} $M = (G, \{P_i\})$ over variables
$V_1,\dots,V_n$ with finite domains $D_1,\dots,D_n$ consists of a DAG $G$ on the
$V_i$ and, for each $i$, a conditional probability table $P_i(v \mid \mathit{pa})$
for $v \in D_i$ and $\mathit{pa} \in \prod_{j \in \PA_i} D_j$. It induces
$\Pr_M(v_1,\dots,v_n) = \prod_i P_i(v_i \mid \mathit{pa}_i)$. $M$ is
\emph{positive} when every table entry is strictly positive. For $X \subseteq V$
and an assignment $x$, the \emph{submodel} $M_x$ replaces $P_i$ by the point
mass at $x_i$ for each $V_i \in X$ and deletes the incoming edges of $X$, and
$\Pr_M(y \mid \dok(x)) := \Pr_{M_x}(y)$.
\end{definition}

A structural causal model refines this by writing each $V_i = f_i(\PA_i, U_i)$
with $U_i$ independent exogenous variables; a CBN is an SCM with the $U_i$
marginalised. The embedding below is of the CBN; \S\ref{ssec:cf} recovers the
$U_i$ as the record of the resolver's draws.

\subsection{The variable gadget}

\begin{definition}[Fan-out]\label{def:fanout}
For a variable $V_i$ with children $V_{j_1},\dots,V_{j_k}$, allocate fresh
channels $v_i^{(j_1)},\dots,v_i^{(j_k)}$, one per child, and a readout channel
$v_i^{\circ}$. Emitting $V_i$'s value $d$ is the parallel composition
$v_i^{(j_1)}!(d) \mid \cdots \mid v_i^{(j_k)}!(d) \mid v_i^{\circ}!(d)$.
\end{definition}

\begin{definition}[The gadget]\label{def:gadget}
For each $V_i$ with parents $V_{p_1},\dots,V_{p_m}$ and domain $D_i =
\{d_1,\dots,d_s\}$, allocate a fresh ancilla channel $w_i$ and write
\begin{lstlisting}
new wi in {
    wi!(d1) | ... | wi!(ds)
  | for( p1 <- vp1_i & ... & pm <- vpm_i
       & u1 <- wi & ... & us <- wi
         where  [ *p1 = pa1 and ... and *pm = pam and *u1 = v ] * P_i(v | pa)
              + ...   // one summand per (pa, v) in the table
       )
    { emit_i( *u1 ) }
}
\end{lstlisting}
where \texttt{vpj\_i} is the copy $v_{p_j}^{(i)}$ of parent $j$ allocated to
this gadget, $\mathit{emit}_i$ is the fan-out of Definition~\ref{def:fanout},
the crisp bracket is the indicator embedding, and the sum has one summand for
each $\mathit{pa} \in \prod_j D_{p_j}$ and $v \in D_i$. A root variable has
$m = 0$.
The \emph{model term} $\sem{M}$ is the parallel composition of the gadgets under
one \textsf{new} binding all the channels.
\end{definition}

The join consumes all $s$ ancilla messages, so there is no residual on $w_i$
and no which-path record; the candidates are the $s!$ bijections from
$u_1,\dots,u_s$ to $D_i$, and the continuation mentions only $u_1$, so the $s!$
candidates fall into $s$ outcome classes of $(s-1)!$ each, one per value of
$u_1$. The clause is constant on each class. Under normalisation within the
contention set the $(s-1)!$ multiplicity cancels, and the probability that the
gadget emits $v$ given parent data $\mathit{pa}$ is $P_i(v \mid \mathit{pa})$.

\begin{lemma}[Staging]\label{lem:staging}
In $\sem{M}$ every receipt is its own contention set, each gadget fires exactly
once, and the causal order of every history is a linear extension of $G$.
\end{lemma}

\begin{proof}
Every channel has exactly one receiver (fan-out gives each child its own copy)
and every ancilla channel is private to its gadget, so no two receipts share a
message or a receipt; contention sets are singletons. A gadget is enabled
exactly when all its parent copies have been emitted, which is
\cite[Prop.~8.3]{GradedV4}: the circuit's DAG is the term's causal order.
\end{proof}

\begin{theorem}[Adequacy of the embedding]\label{thm:adequacy}
Causal histories of $\sem{M}$ reaching a normal form are in bijection with
assignments $(v_1,\dots,v_n) \in \prod_i D_i$, read off the readout channels,
and the likelihood \eqref{eq:lik} of the history for $(v_1,\dots,v_n)$ is
$\Pr_M(v_1,\dots,v_n)$.
\end{theorem}

\begin{proof}
By Lemma~\ref{lem:staging} a history is determined, up to the interleaving the
quotient removes, by the outcome class chosen at each gadget, i.e.\ by the
vector of values; conversely every vector is realised since every ancilla value
is present. The factor contributed by gadget $i$ is its restricted-and-normalised
clause value, $P_i(v_i \mid \mathit{pa}_i)$ by the computation following
Definition~\ref{def:gadget}, and \eqref{eq:lik} multiplies the factors.
\end{proof}

Theorem~\ref{thm:adequacy} says the truncated-factorisation \emph{form} is not
imposed on the semantics; it \emph{is} the semantics, and the staging that makes
each gadget its own contention set is what makes the event structure
confusion-free in the sense of \cite{VaraccaVolzerWinskel}, which is the
condition under which local probabilistic choice composes to a global measure
at all.

\begin{remark}[The Shor gadget]\label{rem:shor}
Definition~\ref{def:gadget} is the one-qubit gadget of \cite[Lemma~8.1]{GradedV4}
and \cite[\S10]{GradedQuantum} --- ancilla channel carrying the domain, a join
that consumes everything, a continuation that re-emits without mentioning what
was consumed, and a clause in disjunctive normal form indexed by (input, output)
--- with three changes: the domain is $D_i$ rather than $\{0,1\}$, the inputs
are several parents rather than one wire, and the matrix the clause realises is
\emph{stochastic} rather than unitary. The resemblance is exact and is a
statement about the semiring. The quantum gadget realises every $2 \times 2$
matrix over $\Cx$ because the clause may mention both the input datum and the
ancilla assignment; the present gadget realises every stochastic matrix over
$\Rnn$ for the same reason, and the restriction to stochastic rather than
arbitrary non-negative matrices is supplied by normalisation within the
contention set rather than by the clause. The two erasures of
\cite[Rem.~8.2]{GradedV4} are both present: the join's totality erases which
ancilla message was which (recombination within a superposition, inert here
since there is no phase to recombine), and the continuation's silence about the
parents erases the input basis (which is what makes the gadget a linear map
into a fixed codomain, i.e.\ a conditional probability table rather than a
family of unrelated tables). What differs is what the two calculi then do with
the matrix: over $\Cx$ one arranges for columns to interfere and reads a
detector that never fires; over $\Rnn$ one reads the trace. Modular
exponentiation in Shor is ordinary rholang running inside every branch at no
cost \cite[Cor.~8.6]{GradedV4}; here the analogous statement is
Proposition~\ref{prop:transparent}, that deterministic subcomputations of a
causal model neither carry probability nor admit intervention.
\end{remark}

\subsection{Surgery}

\begin{theorem}[Graph surgery is restriction of the ancilla race]
\label{thm:surgery}
Let $M$ be positive and $V_i \in V$, $v \in D_i$. Then
\[
  \sem{\,\dok(u_1 = v)\ \ink\ \sem{M}\,} \;=\; \sem{M_{V_i = v}}
\]
as weighted transition systems up to structural congruence, where $u_1$ is the
first ancilla pattern of gadget $i$. Hence
$\Pr_{\sem{M}}(y \mid \dok(V_i = v)) = \Pr_{M}(y \mid \dok(V_i = v))$ for every
$y$.
\end{theorem}

\begin{proof}
At gadget $i$ the restricted clause is the original clause times
$\crisp{*u_1 = v}$, which is $P_i(v \mid \mathit{pa})$ on the $(s-1)!$
candidates in the class $u_1 \mapsto v$ and $0$ elsewhere. The contention set is
the gadget alone (Lemma~\ref{lem:staging}), so normalisation divides by
$(s-1)!\,P_i(v \mid \mathit{pa})$, which is nonzero by positivity, and the class
fires with probability $1$ \emph{for every value of the parents}. The gadget's
normalised dynamics is therefore the point mass at $v$ independent of
$\mathit{pa}$, which is exactly the mechanism of $M_{V_i=v}$ at $V_i$; the
parents are still consumed, so the term is unchanged elsewhere; and every other
gadget is unaffected by Proposition~\ref{prop:lowering}, the model term being
flat. Theorem~\ref{thm:adequacy} applied to $M_{V_i=v}$ gives the distribution.
\end{proof}

The theorem is the reason to address the \emph{ancilla} binding. Restricting a
parent pattern $p_j$ instead would select on inputs --- a conditioning --- not
sever a mechanism. The embedding fixes that the value of a variable is carried
by its ancilla binding, and that is the whole of what makes the surgery exact:
the gadget's own gate is $\top$, so restriction and override coincide
(Observation~\ref{obs:restrict}), and the weights that encoded the mechanism are
the thing normalisation discards.

\begin{remark}[Positivity is starvation]\label{rem:positivity}
If $P_i(v \mid \mathit{pa}) = 0$ for some reachable $\mathit{pa}$, the
restricted gadget has clause $0$ on every candidate when those parents arrive,
and the receipt waits forever: the history deadlocks by starvation, in the sense
of \cite{ProofSearch}. This is exactly where Pearl's $\Pr(y \mid \dok(x))$ is
undefined without positivity, and the embedding renders the undefinedness as a
mortal computation's death rather than as a division by zero. Positivity of
the CPT is decidable by inspection on finite domains; the general question
``is this $\dok$ vacuous or fatal'' is reachability (\S\ref{sec:complexity}).
\end{remark}

\begin{corollary}[The do-calculus is derived]\label{cor:rules}
For positive $M$ and disjoint $X, Y, Z, W \subseteq V$, the three rules of
\cite{Pearl95} hold of $\sem{M}$ under $\dok$ as defined here:
\begin{enumerate}[leftmargin=1.6em,itemsep=1pt]
\item $\Pr(y \mid \dok(x), z, w) = \Pr(y \mid \dok(x), w)$ if
      $(Y \dsep Z \mid X, W)_{G_{\overline{X}}}$;
\item $\Pr(y \mid \dok(x), \dok(z), w) = \Pr(y \mid \dok(x), z, w)$ if
      $(Y \dsep Z \mid X, W)_{G_{\overline{X}\underline{Z}}}$;
\item $\Pr(y \mid \dok(x), \dok(z), w) = \Pr(y \mid \dok(x), w)$ if
      $(Y \dsep Z \mid X, W)_{G_{\overline{X}\,\overline{Z(W)}}}$,
\end{enumerate}
with $Z(W)$ the $Z$-nodes that are not ancestors of $W$ in $G_{\overline{X}}$,
conditioning read as post-selection on histories, and $\dok$ on several
variables as the conjunction $\&$ of their ancilla restrictions.
\end{corollary}

\begin{proof}
By Theorems~\ref{thm:adequacy} and \ref{thm:surgery} and
Proposition~\ref{prop:compose}, every interventional distribution of $\sem{M}$
is the corresponding distribution of a submodel of $M$, and post-selection on
histories is conditioning in that submodel. The rules are Pearl's theorem for
the CBN $M$ \cite{Pearl95,Pearl09}.
\end{proof}

The content of the corollary is not that the rules are true --- Pearl proved
that --- but that the object they are stated for is literally present in the
semantics, so that nothing is being modelled. The d-separation premises are
statements about $G$, and $G$ is the causal order of the staged term
(Lemma~\ref{lem:staging}); we conjecture in \S\ref{sec:open} that they are
expressible as crisp-sort behavioural formulae over $\sem{M}$ and so dischargeable
by the model checker.

\subsection{Counterfactuals}\label{ssec:cf}

Pearl's third rung asks what $Y$ would have been had $X$ been $x$, in the
world that actually occurred. In an SCM this is abduction (infer $U$ from the
evidence), action (mutilate), prediction (compute).

\begin{definition}[Exogenous record]\label{def:record}
A \emph{run} of $\sem{M}$ is a causal history together with, at each gadget, the
candidate that fired. Its \emph{exogenous record} $u$ is the vector of fired
ancilla classes. A run is \emph{replayed under} $\dok(C)$ from a pre-state $S$
by executing $\dok(C)\ \ink\ S$ with the resolver instructed, at each gadget
whose contention set is unchanged from the recorded run, to fire the recorded
class if it is still a candidate of nonzero restricted weight, and otherwise to
sample.
\end{definition}

\begin{theorem}[Counterfactuals are replay]\label{thm:cf}
Let $M$ be positive, let $u$ be the exogenous record of a run with outcome
$(v_1,\dots,v_n)$, and let $S$ be the initial state. Then replaying under
$\dok(u^{(i)}_1 = x)$ from $S$ yields, at every gadget other than $i$ and its
descendants, the recorded value, and at $i$ and its descendants the values
$Y_{V_i = x}(u)$ of the SCM whose exogenous variables are the ancilla draws.
On the staged fragment the ``otherwise sample'' branch of
Definition~\ref{def:record} is never taken.
\end{theorem}

\begin{proof}
Lemma~\ref{lem:staging}: every gadget's contention set is itself, and its
candidate set depends only on which parent values arrived, every ancilla value
being present. So the recorded class is always a candidate; it has nonzero
restricted weight by positivity except at gadget $i$, where the restriction
forces $x$. Non-descendants of $V_i$ receive the same parent values as in the
run and fire the recorded class; descendants receive the mutilated values and
fire their recorded class, which is $f_j(\mathit{pa}_j', u_j)$ for the SCM in
which $u_j$ selects the class --- the standard reduction of a CBN to an SCM by
inverse-CDF sampling, with the class index as $U_j$.
\end{proof}

The ``otherwise sample'' clause is where the general case lives: off the staged
fragment an intervention can change what rests on a channel and thereby which
candidates exist, and ``the same draw'' is no longer well defined. That is the
dynamic-graph obstacle of \S\ref{sec:open} seen from the third rung.

\subsection{The dictionary, completed}

\begin{table}[t]
\centering\small
\begin{tabular}{@{}>{\raggedright\arraybackslash}p{0.27\textwidth}>{\raggedright\arraybackslash}p{0.34\textwidth}>{\raggedright\arraybackslash}p{0.31\textwidth}@{}}
\toprule
\textbf{Pearl} & \textbf{graded \rhoc{} over $\Rnn$} & \textbf{f1r3node}\\
\midrule
variable $V_i$ & gadget $i$; value = fired ancilla class & a COMM event\\
parents $\PA_i$ & data on the channels the pattern and clause read & the contention set, from the pre-state\\
mechanism $P_i(\cdot\mid\mathit{pa})$ & the clause, in DNF & the \texttt{condition} field, pure evaluator\\
exogenous $U_i$ & the resolver's draw at the ancilla race & the executing validator's resolution\\
DAG $G$ & causal order of the staged term & --- \\
$\Pr(y \mid x)$ & post-selection on histories, $1/\Pr(x)$ retries & filtering finalised traces\\
$\Pr(y \mid \dok(x))$ & $\dok(u_1 = x)\ \ink\ \sem{M}$ & elaboration rewrite / membrane\\
$\sigma$-intervention & $\dok(y \models \psi)$ & same\\
$Y_x(u)$ & replay under $\dok$ with the record held & fork from pre-state, replay the trace\\
positivity & no restricted receipt starves & no deadlock in the trace\\
d-separation & behavioural formula over the staged term (conj.) & the Softwhere checker\\
\bottomrule
\end{tabular}
\caption{The dictionary.}\label{tab:dictionary}
\end{table}

Table~\ref{tab:dictionary} collects the correspondence. The right-hand column
is the subject of \S\ref{sec:execution}.

%%%%%%%%%%%%%%%%%%%%%%%%%%%%%%%%%%%%%%%%%%%%%%%%%%%%%%%%%%%%%%%%%%%%%%%%%%%%%%%
\section{Relation to the $\sigma$-calculus}\label{sec:sigma}

Pearl's atomic $\dok$ replaces a mechanism by a constant. Correa and Bareinboim
\cite{CorreaBareinboim} generalise to \emph{$\sigma$-interventions}, which
replace the mechanism of $X$ by an arbitrary new one $P^*(x \mid \mathit{pa}^*)$
--- possibly stochastic, possibly conditional on a different parent set --- and
give a calculus whose rules are stated on the graph $G_\sigma$ in which $X$'s
incoming edges are those of the \emph{new} mechanism. Atomic $\dok$ is the case
$\mathit{pa}^* = \emptyset$ and $P^*$ a point mass; the soft and conditional
interventions of \cite{EberhardtScheines} are intermediate cases.

The construct of \S\ref{sec:construct} is natively a $\sigma$-intervention, and
the three forms it takes have three induced graphs.

\begin{proposition}[Induced graphs]\label{prop:sigmagraph}
Let gadget $i$ of $\sem{M}$ be placed under $\dok(u_1 \models \psi)$. Then the
restricted gadget is a $\sigma$-intervention on $V_i$ whose new mechanism is
$P_i(v \mid \mathit{pa}) \cdot \psi\sem{v, \cdot}$ normalised, and whose parent
set in $G_\sigma$ is
\[
  \PA^*_i \;=\; \PA_i \cup \mathrm{fn}(\psi)\setminus\{u_1\}
\]
when $\psi$ is not an indicator of a singleton, and $\mathrm{fn}(\psi)\setminus\{u_1\}$
when it is. Here $\mathrm{fn}(\psi)$ are the names of variables (readout
channels of other gadgets, or data in scope) that $\psi$'s crisp subformulae
read.
\end{proposition}

\begin{proof}
Restriction multiplies, so the new mechanism depends on $\mathit{pa}$ through
$P_i$ unless the restriction leaves a single class, in which case normalisation
discards $P_i$ (Theorem~\ref{thm:surgery}); and it depends on whatever $\psi$
reads. The parent set of a mechanism is the set of variables it reads.
\end{proof}

Three readings follow. The \emph{set-valued} $\dok(u_1 \in S)$ with $|S| \geq 2$
keeps the original edges: it conditions the mechanism on $S$, which is a
$\sigma$-intervention in which $P^*$ is $P_i$ restricted and renormalised, and
Pearl's three rules do not apply to it verbatim because no edge is deleted. The
\emph{graded} $\dok(u_1 \models \psi)$ with $\psi$ reading only $u_1$ reweights
the mechanism by a function of the value --- a dose, in the vocabulary of
\cite{ExpRegimes}. The \emph{conditional} $\dok(u_1 \models \varphi)$ with
$\varphi$ reading other variables is a policy, a dynamic regime, and its parents
are written in the formula. In every case the graph on which the
$\sigma$-calculus rules are to be evaluated is read off the syntax of the
constraint, which is the same point as the first consequence of
\S\ref{sec:intro}: the parent set is syntactic, now for the intervention as
well as for the model.

\begin{remark}[Override as $\sigma$]
Replacing rather than multiplying --- $\psi_r \mapsto \psi$ --- is a
$\sigma$-intervention with $\PA^* = \mathrm{fn}(\psi)$ outright, and is the
construct one wants when the site's own gate must be severed. It is expressible
as a second mode of the membrane and lowers the same way; we leave it out of
the primitive syntax because on the embedded models it is unnecessary
(Theorem~\ref{thm:surgery}), and off them its complexity is that of
\S\ref{sec:complexity} with no compensating structure.
\end{remark}

\begin{remark}[Cross-site constraints are excluded]
A constraint relating $y_1$ and $y_2$ bound at \emph{different} receipts would
be a correlation across sites: a $\sigma$-intervention with shared exogenous
noise, which is outside both Pearl's and Correa--Bareinboim's calculi and is,
operationally, a join of the two receipts. Definition~\ref{def:do} excludes it by
requiring each $\varphi_i$ to be a constraint on $y_i$ alone; $\&$ is a product of
independent restrictions, which is Pearl's joint $\dok$.
\end{remark}

%%%%%%%%%%%%%%%%%%%%%%%%%%%%%%%%%%%%%%%%%%%%%%%%%%%%%%%%%%%%%%%%%%%%%%%%%%%%%%%
\section{Complexity, and the fragments that tame it}\label{sec:complexity}

The membrane semantics of Definition~\ref{def:membrane} is simple to state and
expensive to run. The expense comes from three independent sources, and a
fourth question --- whether the $\dok$ does anything at all --- is undecidable in
general. We name each, then propose the fragments.

\subsection{Four sources of cost}

\paragraph{The constraint.} $\varphi$ ranges over the full crisp sort, which
contains the modality $\langle K\rangle$ and the fixed point $\nu X$. A
modality-free spatial constraint of depth $d$ costs $O(2^d)$ to evaluate against
a datum (the companion forager charges $\kappa_{\rm sense} = 2^d$ for exactly
this). A modal constraint requires successor exploration to the budget $\beta$,
exponential in $\beta$. A fixed point requires model checking against the
reachable configuration graph, which is PSPACE-hard on finite-state fragments
and undecidable on \rhoc{} at large. And the constraint is evaluated
\emph{per candidate}, since the clause is.

\paragraph{The scope.} Under the membrane semantics the set of constrained
receipts is not known statically: a receipt arriving by reflection must be
checked for a constrained binder when it arrives, and the check is a name
comparison under the hygiene convention, which is cheap, but the
\emph{tracking} of scope through reduction --- which membrane a given receipt is
under, when membranes nest and received code carries its own --- is a run-time
cost proportional to the depth of nesting, paid on every COMM.

\paragraph{The contention set.} Restriction removes candidates and so changes
the normalisation of the contention set they belonged to. Where contention sets
are singletons (Lemma~\ref{lem:staging}) this is a division. Where they are
not, the normaliser is a sum over maximal matchings, which is permanent-shaped
and $\#\mathsf{P}$-hard in general \cite{GradedQuantum}; a $\dok$ that touches a
large contention set forces that sum to be recomputed.

\paragraph{Vacuity and fatality.} Whether any receipt binding $y$ ever fires
inside $P$ is reachability, undecidable in general. So whether a given $\dok$
is vacuous (no constrained race ever arises), effective (some arise and some
candidate always satisfies $\varphi$), or fatal (a constrained race arises with
no satisfying candidate, and starves) cannot be decided for arbitrary $P$. This
is the halting argument of \S\ref{sec:construct} returning as a cost: the
programmer cannot in general know what an intervention will cost before paying
for it.

\subsection{Fragments of the constraint language}

\begin{definition}[Constraint fragments]\label{def:fragments}
\begin{description}[leftmargin=1.6em,itemsep=2pt]
\item[$\mathcal{F}_0$ (values).] $y = Q$ and $y \in S$ for a finite set $S$ of
closed terms; and $\bigoplus_i v_i \tensorv \crisp{y = Q_i}$. Evaluation is a
structural-congruence test against $|S|$ terms: constant in the term, linear in
the datum. This is the Pearl fragment: Theorems~\ref{thm:surgery} and
\ref{thm:cf} and Corollary~\ref{cor:rules} use nothing else.
\item[$\mathcal{F}_1$ (spatial).] Modality-free, fixed-point-free crisp
formulae over the bound datum and the data in scope, of bounded depth $d$; and
graded formulae built from them by $\tensorv$, $\sumv$ and constants. Evaluation
is $O(2^d)$ per candidate and needs no successor exploration; the graded members
are posynomials in their constants \cite[Prop.~5.2]{GradedPLN}, so a learner
can differentiate them. Conditional $\sigma$-interventions
(Proposition~\ref{prop:sigmagraph}) live here.
\item[$\mathcal{F}_2$ (budgeted modal).] $\mathcal{F}_1$ closed under
$\langle K\rangle_\V$ to a declared depth $\beta$. Evaluation explores
successors to depth $\beta$ and is charged the inference budget already present
in the calculus \cite[Rem.~6.5]{GradedV4}. Interventions conditional on what a
site \emph{could} do next --- ``restrict to candidates after which $\varphi$ is
reachable in two steps'' --- live here.
\item[$\mathcal{F}_3$ (fixed points).] The full crisp sort. Not evaluable at
the site in general; admissible in \emph{specification} position, where the
model checker discharges it offline on a finite-state fragment, and then
lowered to $\mathcal{F}_1$ by replacing the fixed point with the finite
disjunction it unfolds to on that fragment.
\end{description}
\end{definition}

\begin{principle}[Where each fragment runs]\label{prin:fragments}
$\mathcal{F}_0$ and $\mathcal{F}_1$ run at the site, in the clause evaluator,
at the cost of a guard. $\mathcal{F}_2$ runs at the site under the inference
budget and is priced by it. $\mathcal{F}_3$ belongs to the checker, not the
clause: it is what one \emph{verifies} about an intervention, not what one
\emph{executes} as one. A $\dok$ whose constraint is in $\mathcal{F}_3$ is a
specification awaiting lowering, and Remark~\ref{rem:specguard} says that
lowering is the construct's whole job.
\end{principle}

\subsection{Type-based constraints on the scope}

The constraint fragments bound the cost of one evaluation. The scope
disciplines bound how many evaluations there are, whether the scope is static,
and whether the contention sets are small. They are types in the sense the
programme already uses: spatial-behavioural formulae imposed on $P$ as
admission conditions, decidable on the fragment and checked once.

\begin{definition}[Scope disciplines]\label{def:types}
\begin{description}[leftmargin=1.6em,itemsep=2pt]
\item[$\mathcal{T}_1$ (staged).] Every channel in scope has exactly one
receiver, every receipt consumes every message on each of its channels, and
data flow uses the linear binding arrow $\lolli$ of \cite{GradedQuantum}
outside continuations' fan-out. Consequence: contention sets are singletons,
the normaliser is a division, the causal order is a static DAG
(Lemma~\ref{lem:staging}), and each gadget fires once.
\item[$\mathcal{T}_2$ (flat).] No drop $*z$ in scope at a position that can
come to contain a receipt; equivalently the stratified \where{} of
\cite{ChoiceFormulae} with ceiling $\bot$ on binders. Consequence: the set of
constrained receipts is static and lowering is exact
(Proposition~\ref{prop:lowering}); no scope tracking at run time.
\item[$\mathcal{T}_3$ (finite-domain).] Every channel a constrained variable is
bound on carries a DDL type with finitely many closed values. Consequence:
$\mathcal{F}_0$ constraints are sets of enumerable size; vacuity and fatality
are decidable by inspecting which values the clause assigns positive weight,
which is the positivity check of Remark~\ref{rem:positivity}.
\item[$\mathcal{T}_4$ (confined).] Every constrained receipt lies within the
deployer's seal in the sense of \cite{CostMonad}, so that the $\dok$ reaches
only code its author pays for. Consequence: the membrane coincides with the
seal tree, the capability to intervene is the capability to deploy, and
Proposition~\ref{prop:nocontext} is not a limitation but the boundary of the
namespace.
\end{description}
\end{definition}

\begin{proposition}[The practical fragment]\label{prop:practical}
Under $\mathcal{T}_1$--$\mathcal{T}_4$ with constraints in $\mathcal{F}_0$:
(i) $\dok(C)\ \ink\ P$ lowers at elaboration and the membrane is erased;
(ii) each constrained candidate costs one structural-congruence test against
$|S|$ closed terms; (iii) vacuity and fatality are decidable; (iv) the induced
graph $G_\sigma$ is the static DAG with the constrained variables' incoming edges
deleted; and (v) the semantics is that of Definition~\ref{def:cbn} exactly.
Under $\mathcal{T}_1$--$\mathcal{T}_3$ with constraints in $\mathcal{F}_1$,
(i)--(iii) hold, $G_\sigma$ is the static DAG with $\PA^*_i$ as in
Proposition~\ref{prop:sigmagraph}, and the semantics is a $\sigma$-intervention
on a CBN.
\end{proposition}

\begin{proof}
(i) is Proposition~\ref{prop:lowering}, whose hypothesis is $\mathcal{T}_2$.
(ii) is the definition of $\mathcal{F}_0$. (iii): under $\mathcal{T}_1$ every
constrained receipt fires exactly once and its candidate set is determined by
its parents' values, which range over a finite product by $\mathcal{T}_3$;
enumerate. (iv) and (v) are Theorem~\ref{thm:surgery} and
Proposition~\ref{prop:sigmagraph}, whose hypotheses are supplied by
$\mathcal{T}_1$ and $\mathcal{T}_3$. The second statement replaces
Theorem~\ref{thm:surgery} by Proposition~\ref{prop:sigmagraph}.
\end{proof}

\begin{table}[t]
\centering\footnotesize
\begin{tabular}{@{}>{\raggedright\arraybackslash}p{0.15\textwidth}>{\raggedright\arraybackslash}p{0.26\textwidth}>{\raggedright\arraybackslash}p{0.19\textwidth}>{\raggedright\arraybackslash}p{0.29\textwidth}@{}}
\toprule
\textbf{constraint} & \textbf{per-candidate cost} & \textbf{needs} & \textbf{supports}\\
\midrule
$\mathcal{F}_0$ & $O(|S|)$ congruence tests & nothing & Pearl's $\dok$, rules, counterfactuals\\
$\mathcal{F}_1$ & $O(2^d)$ & bounded depth $d$ & conditional $\sigma$, doses, policies\\
$\mathcal{F}_2$ & $O(b^\beta)$, $b$ branching & inference budget $\beta$ & interventions on reachable futures\\
$\mathcal{F}_3$ & checker & finite-state fragment & specification only; lower to $\mathcal{F}_1$\\
\midrule
\textbf{scope type} & \multicolumn{3}{l@{}}{\textbf{what it buys}}\\
\midrule
$\mathcal{T}_1$ staged & \multicolumn{3}{>{\raggedright\arraybackslash}p{0.74\textwidth}@{}}{singleton contention sets; static DAG; one firing per gadget}\\
$\mathcal{T}_2$ flat & \multicolumn{3}{>{\raggedright\arraybackslash}p{0.74\textwidth}@{}}{exact lowering; no run-time scope tracking}\\
$\mathcal{T}_3$ finite-domain & \multicolumn{3}{>{\raggedright\arraybackslash}p{0.74\textwidth}@{}}{decidable vacuity and fatality}\\
$\mathcal{T}_4$ confined & \multicolumn{3}{>{\raggedright\arraybackslash}p{0.74\textwidth}@{}}{$\dok$ reaches only paid-for code; the capability to intervene is the capability to deploy}\\
\bottomrule
\end{tabular}
\caption{Fragments and disciplines.}\label{tab:fragments}
\end{table}

The disciplines are not new to the programme. $\mathcal{T}_1$ is the safe
fragment of \cite[\S9]{GradedQuantum} in which $\sem{P}$ is a quantum state;
$\mathcal{T}_2$ is the flat fragment of \cite{ExpRegimes} and the stratified one
of \cite{ChoiceFormulae}; $\mathcal{T}_3$ is what the DDL supplies for free; and
$\mathcal{T}_4$ is the seal discipline of the cost-accounting work. That four
notes written for four purposes draw the same line is the best evidence we have
that the line is real. The route to larger fragments is the one the
monadic-foreach note already proposes: typed calculi translated into the
language with types as spatial-behavioural formulae, System~F first and then
Abramsky's CLL, each of which supplies a terminating fragment on which
$\mathcal{F}_3$ is decidable.

\begin{remark}[What is lost outside the fragment]
Outside $\mathcal{T}_1$ the contention sets are the branching cells of
\cite{AbbesBenveniste} and a $\dok$ that touches one must be understood as a
restriction of the cell's distribution; outside $\mathcal{T}_2$ a $\dok$ is a
genuine run-time membrane; outside $\mathcal{T}_3$ positivity is a reachability
question; outside $\mathcal{T}_4$ an intervention on another's code is a soft
intervention by Proposition~\ref{prop:nocontext}, an experiment rather than a
surgery. None of these is a failure of the construct. They are where Pearl's
calculus stops and the calculus of this programme continues.
\end{remark}

%%%%%%%%%%%%%%%%%%%%%%%%%%%%%%%%%%%%%%%%%%%%%%%%%%%%%%%%%%%%%%%%%%%%%%%%%%%%%%%
\section{Execution}\label{sec:execution}

The construct was designed against the execution framework rather than for a
simulator, and most of what it needs exists.

\paragraph{Lowering.} The normaliser already carries an elaboration pass that
puts every deployment into normative forms so that the reducer's seal record is
read off the term. Lowering $\dok$ under $\mathcal{T}_2$ is one more rewrite in
that pass, and it leaves no run-time trace: the constrained clauses are
ordinary clauses with one more conjunct, evaluated by the pure
\texttt{condition} evaluator the matcher already calls per candidate.

\paragraph{The membrane.} Off $\mathcal{T}_2$ the $\dok$ is a run-time membrane,
and it is the twin of the cost membrane $\{\%\ P\ \%\}[s]$: a scoped decoration
on the cut, read off the term by the reducer, altering the COMM rule within.
The seal tree rooted at the deployer's wallet becomes a seal-and-constraint
tree; a receipt arriving by reflection inherits the innermost constraint as it
inherits the innermost seal. The two decorations commute
(Proposition~\ref{prop:compose} extends to them, since metering does not read
the clause and the clause does not read the meter) and so may be nested in
either order.

\paragraph{Pricing.} An intervention is one COMM on the port, by
Principle~\ref{prin:metered}, plus the marginal cost of the extra conjunct per
candidate evaluation, which under $\mathcal{F}_0$ is a congruence test. The
charge for an $\mathcal{F}_2$ constraint is the inference budget it declares.
Fixing these against the ledger is Obligation~\ref{ob:price}.

\paragraph{Observation, and who does it.} The validators' execution model ---
one node executes and produces a trace, the others replay it to confirm the
state change --- is, in the vocabulary of \cite{Inhabiting}, an
experimenter/spectator split. The executing validator resolves the races and
so performs rung two; the replayers observe and are rung one. $\Pr(y \mid x)$ is
estimated from the finalised trace; it costs nothing beyond storage and is
exactly what \cite[Prop.~4.1]{Inhabiting} says it is: the trace distribution,
and no more.

\paragraph{Counterfactuals.} Theorem~\ref{thm:cf} needs a pre-state and a
record of fired candidates. The order-then-execute design runs each envelope
from the block pre-state and records the winners; the losing candidates are
recomputable from the pre-state deterministically, so nothing further need be
stored. The same record is what the gradient of \cite[Cor.~5.7]{GradedPLN}
accumulates --- one number per firing from the candidates the site already
evaluated --- so one instrumentation serves learning and counterfactual replay.
The stable site identity the replay requires is the quotation of the receipt,
exactly as the heralds of \cite[\S9]{GradedV4} name their sites.

\paragraph{The kernel.} $\Pr(y \mid \dok(x))$ is $e_y^{\!\top}\, T_{\dok(x)}^{\beta}
\,\lvert P\rangle$ for the transfer operator with one column block replaced, and
the adjoint of \cite[Thm.~5.6]{GradedPLN} is the same Neumann series. A batch
of interventional queries and the modal gradient share one sparse matrix-vector
path on the GPU.

%%%%%%%%%%%%%%%%%%%%%%%%%%%%%%%%%%%%%%%%%%%%%%%%%%%%%%%%%%%%%%%%%%%%%%%%%%%%%%%
\section{Inhabiting choice, discharged}\label{sec:inhabiting}

The companion note \cite{Inhabiting} makes one observation three times: a
system whose witnesses reduce confluently can \emph{state} a non-confluent
phenomenon but cannot \emph{inhabit} it, so its identity on witnesses is forced
to be linear-time, and the logic adequate for that identity is forced to be
determinate. In its reading for artificial intelligence it aligns the
linear/branching/causal spectrum with Pearl's ladder --- association,
intervention, counterfactual --- and names the remedy as witnesses that can be
races and an identity on them at least as fine as step bisimulation. The
alignment is by argument. The present note is its constructive content, and we
propose the mechanism as the tool for each of the obligations that note leaves
owed.

\paragraph{Rung two is a graded race under a $\dok$.} The note's rung two is
``what an experimenter supplying contexts can learn,'' characterised by
context-labelled bisimulation. Proposition~\ref{prop:nocontext} sharpens this:
a context supplies apertures and cannot force them, so the experimenter's power
over $\Rnn$ is exactly the $\dok$ membrane on the code it owns together with
soft interventions on the rest. Theorem~\ref{thm:surgery} says that on the
staged fragment that power is Pearl's rung two on the nose. Rung three is
Theorem~\ref{thm:cf}: replay with the record held.

\paragraph{The vending-machine benchmark becomes a program.} The smallest
behaviour invisible to a trace is Milner's pair
\[
  P = \mathit{coin}.(\mathit{tea} + \mathit{coffee}), \qquad
  Q = \mathit{coin}.\mathit{tea} + \mathit{coin}.\mathit{coffee},
\]
and the observable a trace cannot recover is a \emph{refusal}: coin in, tea
requested, nothing out. Render the machines with races in the usual way, with
$y$ the variable bound at the drink-selecting receipt. Then
\[
  \dok(y = \mathit{tea})\ \ink\ (\mathit{coin}!() \mid P)
  \quad\text{and}\quad
  \dok(y = \mathit{tea})\ \ink\ (\mathit{coin}!() \mid Q)
\]
differ in exactly one way: the first always delivers tea, and the second, on
the branch where $Q$ has committed to coffee before the coin, has a restricted
receipt with no satisfying candidate and \emph{starves}. The refusal that never
appears in a log of completed runs is a deadlock in the intervened term. The
benchmark the note proposes --- logs from one machine chosen by a hidden coin,
completed actions only, then the refusal question --- is answered by one
$\dok$ and one observation of whether anything was delivered, which is the
``one attempted tea after one coin'' the note says settles the question; here it
is a term one can run.

\paragraph{The audit as a theorem.} The note's verification reading says a
$\lambda$-term prover can certify a concurrent system only relative to a
stipulated model of its non-determinism, and that the faithfulness of the model
is not auditable by any means the prover possesses. Its obligation is to say
which class of tests such a prover can run and prove the class contained in
trace observation. Proposition~\ref{prop:nocontext} is half of that statement:
a test that cannot force a race is a context, and a context's observations are
contained in the support reading. The other half is that a $\dok$ is precisely
a test outside that class, and that on a process runtime the certificate can
contain one.

\paragraph{The three requirements as a theorem.} The note's closing section
argues that recursive self-improvement needs an improvement operator that acts
on witnesses that are races, preserves branching through evaluation, and funds
the branches that lose. A $\dok$ is an operator on witnesses that changes their
branching --- it is the only kind there is, by \S\ref{sec:carrier} --- and the
set-valued form $\dok(y \in S)$ with $S$ the \emph{losing} candidates is the
operator that funds them: it restricts the race away from its current winner.
That an improvement operator definable on a determinate language cannot increase
branching is Proposition~\ref{prop:transparent} read for the whole language;
that one definable on \rhoc{} with the cost endofunctor can is this note.

\paragraph{Hedges that are inherited.} The note's last row asks what leaves the
saddle point: tactics as witnesses, certificates that run, hedges that are
inherited. A hypothesis of the mortal scientist is a clause; a clause under a
$\dok$ is a hypothesis together with the experiment that would test it; and a
germ carrying both is a hedge that is inherited. Which is to say that the
construct of this note is what the scientist's germs were missing in order to
be rung-two objects rather than rung-one ones, and the reproduction note's
crossover on quoted code applies to it unchanged, since a $\dok$ is a term.

%%%%%%%%%%%%%%%%%%%%%%%%%%%%%%%%%%%%%%%%%%%%%%%%%%%%%%%%%%%%%%%%%%%%%%%%%%%%%%%
\section{Open problems}\label{sec:open}

\begin{obligation}[d-separation as a behavioural formula]\label{ob:dsep}
Corollary~\ref{cor:rules} evaluates the d-separation premises on $G$. Since
$G$ is the causal order of the staged term, the premises should be crisp-sort
formulae over $\sem{M}$ --- ``no history has a causal path from a $Z$-firing to a
$Y$-firing not passing through $X$ or $W$'' is a $\nu$-property --- dischargeable
by the model checker. Stating and proving the correspondence would make the
rules' side conditions checkable in the same instrument as their conclusions.
\end{obligation}

\begin{obligation}[The rules off the staged fragment]\label{ob:general}
Outside $\mathcal{T}_1$ the contention sets are branching cells and the event
structure may be confused. The conjecture is that the $\sigma$-calculus rules
hold on the confusion-free fragment with d-separation read off the event
structure, and that on the confused fragment they hold per branching cell with
the cell's distribution in place of the mechanism. The results of
\cite{VaraccaVolzerWinskel,AbbesBenveniste} are the starting point and should
be verified against the sources before being leaned on.
\end{obligation}

\begin{obligation}[The price of a $\dok$]\label{ob:price}
Principle~\ref{prin:metered} says an intervention is a COMM. The marginal
charge of the constraint conjunct, the budget charge of an $\mathcal{F}_2$
constraint, and the charge for a counterfactual replay have not been fixed
against the ledger. Until they are, ``an intervention costs one phlo'' names an
intention.
\end{obligation}

\begin{obligation}[Identifiability across namespaces]\label{ob:ident}
Pearl's identifiability question --- when is $\Pr(y \mid \dok(x))$ computable
from observational data in the presence of hidden variables --- has a sharp
rendering here: the hidden variables are the code one cannot read, the boundary
is the namespace, and the observational data is the shard's finalised trace.
The ID algorithm of \cite{ShpitserPearl} should transfer to the staged fragment
directly; what it says about a shard is the question.
\end{obligation}

\begin{obligation}[Continuous domains]\label{ob:continuous}
The gadget enumerates the domain on the ancilla channel. A continuous variable
needs the DDL's numeric types to carry an ancilla as a density, and the
contention-set normalisation to become an integral. The Gillespie reading of
\cite{WeightedGSLT} is the natural home.
\end{obligation}

\begin{obligation}[Cycles]\label{ob:cycles}
A recursive term has an acyclic event-level graph and a cyclic variable-level
one. The cyclic SCMs of \cite{BongersEtAl} and the graded fixed point of
\cite{GradedPLN} should meet here, with the budget $\beta$ as the unrolling
depth.
\end{obligation}

%%%%%%%%%%%%%%%%%%%%%%%%%%%%%%%%%%%%%%%%%%%%%%%%%%%%%%%%%%%%%%%%%%%%%%%%%%%%%%%
\section{Conclusion}\label{sec:conc}

A graded program over the non-negative reals was a causal model before anyone
said so: its likelihood factorises over sites, its parents are written in its
clauses, and its exogenous noise is the resolver the programme had already made
a parameter. What it lacked was an operator that forces a race, and the operator
turns out to want a specific home. It cannot sit at a node, because whether a
node will ever occur is undecidable; it sits at the cut, as a policy conditional
on a race arising, and it lowers into the where-clause by the normal-form
theorem for weight maps, so it adds nothing to the calculus that was not already
expressible --- only a place to say it.

With it, Pearl's calculus embeds exactly, and the embedding is the Shor gadget
with a stochastic matrix in place of a unitary one. Intervening on the ancilla
binding deletes incoming edges because normalisation discards the weights that
encoded the mechanism; positivity is a receipt that does not starve; the three
rules are theorems; counterfactuals are replay. Set-valued and graded
constraints land in the $\sigma$-calculus with the induced graph read off the
free names of the constraint. The construct is expensive to run naively, in
ways that are each separately familiar from the programme's other notes, and
the fragments that tame it are the fragments those notes had already drawn.

The mechanism is offered as the constructive half of the inhabiting-choice
argument. A race under a $\dok$ is rung two; its replay is rung three; and the
refusal that no corpus of completed runs can reveal is, in the intervened
term, a receipt that waits for a tea that will never come.

%%%%%%%%%%%%%%%%%%%%%%%%%%%%%%%%%%%%%%%%%%%%%%%%%%%%%%%%%%%%%%%%%%%%%%%%%%%%%%%
\begin{thebibliography}{99}\raggedright

\bibitem{Pearl95}
J.~Pearl.
\newblock Causal diagrams for empirical research.
\newblock \emph{Biometrika} 82(4):669--688, 1995.

\bibitem{Pearl09}
J.~Pearl.
\newblock \emph{Causality: Models, Reasoning, and Inference}.
\newblock Cambridge University Press, second edition, 2009.

\bibitem{Pearl93}
J.~Pearl.
\newblock Graphical models, causality and intervention.
\newblock \emph{Statistical Science} 8(3):266--269, 1993.

\bibitem{BareinboimHierarchy}
E.~Bareinboim, J.~D.~Correa, D.~Ibeling, and T.~Icard.
\newblock On Pearl's hierarchy and the foundations of causal inference.
\newblock In \emph{Probabilistic and Causal Inference: The Works of Judea
Pearl}, ACM Books, 2022.

\bibitem{CorreaBareinboim}
J.~D.~Correa and E.~Bareinboim.
\newblock A calculus for stochastic interventions: causal effect identification
and surrogate experiments.
\newblock In \emph{AAAI 2020}, pp.~10093--10100.

\bibitem{ShpitserPearl}
I.~Shpitser and J.~Pearl.
\newblock Identification of joint interventional distributions in recursive
semi-Markovian causal models.
\newblock In \emph{AAAI 2006}, pp.~1219--1226.

\bibitem{EberhardtScheines}
F.~Eberhardt and R.~Scheines.
\newblock Interventions and causal inference.
\newblock \emph{Philosophy of Science} 74(5):981--995, 2007.

\bibitem{BongersEtAl}
S.~Bongers, P.~Forr\'e, J.~Peters, and J.~M.~Mooij.
\newblock Foundations of structural causal models with cycles and latent
variables.
\newblock \emph{Annals of Statistics} 49(5):2885--2915, 2021.

\bibitem{VaraccaVolzerWinskel}
D.~Varacca, H.~V\"olzer, and G.~Winskel.
\newblock Probabilistic event structures and domains.
\newblock \emph{Theoretical Computer Science} 358(2--3):173--199, 2006.

\bibitem{AbbesBenveniste}
S.~Abbes and A.~Benveniste.
\newblock True-concurrency probabilistic models: branching cells and
distributed probabilities for event structures.
\newblock \emph{Information and Computation} 204(2):231--274, 2006.

\bibitem{CostaShrapnel}
F.~Costa and S.~Shrapnel.
\newblock Quantum causal modelling.
\newblock \emph{New Journal of Physics} 18:063032, 2016.

\bibitem{AllenEtAl}
J.-M.~A.~Allen, J.~Barrett, D.~C.~Horsman, C.~M.~Lee, and R.~W.~Spekkens.
\newblock Quantum common causes and quantum causal models.
\newblock \emph{Physical Review X} 7:031021, 2017.

\bibitem{LarsenSkou}
K.~G.~Larsen and A.~Skou.
\newblock Bisimulation through probabilistic testing.
\newblock \emph{Information and Computation} 94(1):1--28, 1991.

\bibitem{MilnerCCS}
R.~Milner.
\newblock \emph{Communication and Concurrency}.
\newblock Prentice Hall, 1989.

\bibitem{MeredithRadestock}
L.~G.~Meredith and M.~Radestock.
\newblock A reflective higher-order calculus.
\newblock \emph{Electronic Notes in Theoretical Computer Science}
141(5):49--67, 2005.

\bibitem{GradedV4}
L.~G.~Meredith.
\newblock Graded where-clauses: non-crisp truth values as the syntax of
non-determinism resolution in the rho calculus, version 4.
\newblock F1R3FLY.io working note, \texttt{fuzzyware/graded-where-v4}, August
2026.

\bibitem{GradedQuantum}
L.~G.~Meredith.
\newblock Races decided by the Born rule: the quantum instance of the graded
where-clause, version 6.
\newblock F1R3FLY.io working note, \texttt{fuzzyware/graded-where-quantum-v6},
August 2026.

\bibitem{GradedPLN}
L.~G.~Meredith.
\newblock Graded where-clauses over the reals: one-dimensional grading as the
mortal scientist's gradient, and probabilistic logic as the metric that
gradient requires.
\newblock F1R3FLY.io working note, \texttt{fuzzyware/graded-where-pln}, August
2026.

\bibitem{Omnibus}
L.~G.~Meredith.
\newblock Observation disciplines: how logics are generated from term
structure, and why graded where-clauses are the mortal scientist's instrument.
\newblock F1R3FLY.io working note, \texttt{rho-mind/observation-disciplines},
August 2026.

\bibitem{ChoiceFormulae}
L.~G.~Meredith.
\newblock Choice principles as graded where-clauses: one programming paradigm
at every altitude of the Weihrauch lattice.
\newblock F1R3FLY.io working note, \texttt{choice-types/choice-formulae},
August 2026.

\bibitem{ExpRegimes}
L.~G.~Meredith.
\newblock Experimental regimes: laboratories for universal negatives.
\newblock F1R3FLY.io working note, \texttt{pnp-experiments/experimental-regimes},
2026.

\bibitem{Inhabiting}
L.~G.~Meredith.
\newblock Inhabiting choice: what a confluent witness cannot be.
\newblock F1R3FLY.io working note,
\texttt{proof-dynamics-and-reasoning/inhabiting-choice}, November 2026.

\bibitem{ProofSearch}
L.~G.~Meredith.
\newblock Proof search as concurrent reaction.
\newblock F1R3FLY.io working note,
\texttt{proof-dynamics-and-reasoning/proof-search}, 2026.

\bibitem{CostMonad}
L.~G.~Meredith.
\newblock Cost accounting as an endofunctor on continued interactive GSLTs.
\newblock F1R3FLY.io working note,
\texttt{cost-accounting-as-monad/continued-gslt-cost-v2}, 2026.

\bibitem{WeightedGSLT}
L.~G.~Meredith.
\newblock Weighted GSLTs: stochastic and quantum simulation for decorated
continued interactive GSLTs.
\newblock F1R3FLY.io working note, \texttt{rho-life/weighted-gslt}, 2026.

\bibitem{Gillespie}
D.~T.~Gillespie.
\newblock Exact stochastic simulation of coupled chemical reactions.
\newblock \emph{Journal of Physical Chemistry} 81(25):2340--2361, 1977.

\bibitem{HennessyMilner}
M.~Hennessy and R.~Milner.
\newblock Algebraic laws for nondeterminism and concurrency.
\newblock \emph{Journal of the ACM} 32(1):137--161, 1985.

\end{thebibliography}

\end{document}
